\documentclass[12pt]{article}

% arXiv categories: cs.AI cs.CL cs.LG
\usepackage{amsmath, amssymb, amsthm}
\usepackage{geometry}
\usepackage{cite}
\usepackage{enumitem}
\usepackage{hyperref}
\usepackage{booktabs}

\geometry{margin=1in}

\newtheorem{definition}{Definition}
\newtheorem{proposition}{Proposition}
\newtheorem{remark}{Remark}
\newtheorem{corollary}{Corollary}

\title{Hallucination as Frame Misactivation:\\
A Structural Account of LLM Error Under the\\
Interpretive Stack Model}

\author{Minamo Minamoto\\
\small Independent\\
\small \texttt{minamominamoto4f5683f6@gmail.com}\\
\small \texttt{https://orcid.org/0009-0002-1201-5704}}

\date{April 2026}

\begin{document}

\maketitle

\begin{abstract}
Large language model (LLM) hallucination is standardly explained as
a failure of uncertainty calibration: the model generates confident
output in regions of low epistemic warrant. We argue that this
account, while partially correct, is structurally incomplete. Many
hallucination types---including contextual confabulation, role
confusion, narrative overcompletion, and cross-domain contamination
---are not primarily failures of confidence calibration but failures
of \emph{frame selection}: the model activates an interpretive frame
that is inappropriate for the current input context, and generates
fluent, internally coherent output within that frame.

We propose a structural account of hallucination grounded in an
interpretive stack model \cite{minamoto2026accumulation}: an LLM implicitly maintains a layered set of
interpretive frames acquired during pretraining and fine-tuning,
and hallucination occurs when the \emph{selection rule} activates
the wrong frame in response to the current input. Under this model,
the apparent diversity of hallucination types can be unified within
a single explanatory framework with a predictable structure.

This account has three consequences. First, it generates a
classification of hallucination types by their frame-selection
failure mechanism, unifying phenomena that existing taxonomies treat
as distinct. Second, it reframes retrieval-augmented generation (RAG),
chain-of-thought (CoT) prompting, and system prompts as practical
interventions on frame activation rather than knowledge augmentation
or reasoning support. Third, it identifies \emph{Selection Stability
Capacity} (SSC)---the ability to maintain the correct frame across
diverse input contexts---as the structural property that debiasing
interventions should target, and explains why increasing training
data volume alone is insufficient to achieve it.

We derive four testable predictions distinguishing the frame-selection
account from the uncertainty-calibration account, and propose
mechanistic interpretability as the primary empirical route for
verification. This framework does not merely reinterpret hallucination:
it predicts systematic shifts in error type distribution under different
prompting conditions, retrieval strategies, and fine-tuning regimes---
predictions that are falsifiable and distinct from those of the
calibration account.
\end{abstract}

%-------------------------------------------------------------------
\section{Introduction}
%-------------------------------------------------------------------

Hallucination in large language models---the generation of fluent,
confidently-stated content that is factually incorrect, contextually
inappropriate, or unsupported by the input---has become a central
problem in AI deployment. The standard account attributes
hallucination to \emph{epistemic failure}: the model lacks reliable
uncertainty estimates and generates high-confidence output in
low-confidence regions \cite{kadavath2022, lin2022, xiong2024}.
The canonical prescription follows directly: improve calibration
through temperature scaling, uncertainty-aware training, or
reinforcement learning from human feedback (RLHF).

We do not dispute that calibration failures contribute to
hallucination. We argue, however, that calibration failure is
insufficient to explain a significant class of hallucination
phenomena. Consider the following:

\begin{itemize}
\item A model asked about a historical figure invents plausible-
sounding but false biographical details, even when it could in
principle express uncertainty.
\item A model playing a specific role (e.g., a medical advisor)
generates output appropriate to a different role (e.g., a general
conversationalist) when the input shifts register slightly.
\item A model given a document to summarize introduces content
from a related but distinct document not in context.
\item A model generating a narrative continues the narrative in a
direction consistent with its training distribution rather than with
the explicit constraints given in the prompt.
\end{itemize}

In each of these cases, the hallucinated content is not random noise
but \emph{coherent within a frame}---the model is generating output
that would be appropriate if a different context, role, or knowledge
base were active. The error is not that the model is uncertain; it
is that the model is confidently operating within the wrong
interpretive frame.

We propose that this class of hallucinations is better understood as
\emph{frame misactivation}: the selection rule that governs which
interpretive frame is active has chosen an inappropriate frame for
the current input, and the model generates fluent output within that
frame without detecting the mismatch. This is structurally distinct
from uncertainty-calibration failure and calls for a different
diagnostic and remediation approach.

The present paper develops this account within the interpretive stack
framework of \cite{minamoto2026accumulation}, applies
it to derive a classification of hallucination types, reinterprets
RAG and CoT as frame-selection interventions, and derives testable
predictions that distinguish the frame-selection account from the
calibration account.

%-------------------------------------------------------------------
\section{The Interpretive Stack Model}
%-------------------------------------------------------------------

We introduce the four constructs required for the hallucination
analysis. The framework is developed formally in companion work
\cite{minamoto2026accumulation}; the present paper
is self-contained with respect to the definitions below.

\begin{definition}[Interpretive Frame]
An \emph{interpretive frame} $F$ is a function from inputs to outputs
induced by a specific configuration of a language model's weight
space. For an LLM, a frame corresponds to a computational subgraph
that governs generation for a class of inputs---for example, the
``medical knowledge'' frame, the ``narrative continuation'' frame,
or the ``instruction following'' frame.

\textbf{Operational grounding.} We do not assume that frames
correspond to cleanly separable weight subsets. Rather, we adopt the
operational definition from mechanistic interpretability research
\cite{conmy2023towards, bricken2023monosemanticity}: a frame is
a cluster of circuits (in the sense of \cite{elhage2021mathematical})
whose causal contribution to output is coherent across a class of
inputs. Operationally, two inputs $e_1$ and $e_2$ activate the
\emph{same frame} if activation patching from $e_1$ to $e_2$ (or
vice versa) leaves output behavior qualitatively unchanged for the
relevant task, and activates a \emph{different frame} if patching
produces a qualitative shift in output type or style. This definition
connects the interpretive stack model to empirically tractable methods
and makes the existence of distinct frames a testable claim rather
than a posit.
\end{definition}

\begin{remark}[Computational-level status of the frame definition]
The activation patching criterion above is a \emph{necessary condition}
for same-frame classification, not a standalone decision procedure.
In practice, applying activation patching requires a prior partition of
inputs into candidate frame classes; the criterion confirms or
disconfirms a proposed classification, not generates one from scratch.
The definition is therefore \emph{empirically tractable given a frame
hypothesis}. The computational form of the selection rule
$\mathcal{S}$---which specific attention heads, MLP layers, or residual
stream directions implement context-to-frame mapping---remains open
(OP2 below). The interpretive stack model is a \emph{computational-level}
theory in Marr's sense: it specifies what the system computes
(frame selection) without yet fully specifying the implementing mechanism.
This is a deliberate starting point, not a limitation.
\end{remark}

\begin{definition}[Interpretive Stack]
An \emph{interpretive stack} $\mathcal{L} = (F_1, F_2, \ldots, F_n)$
is an ordered sequence of interpretive frames implicitly encoded in
the model's weight space. $F_1$ is the most recently installed or
highest-priority frame. Earlier frames are residual structures from
pretraining or prior fine-tuning stages.
\end{definition}

\begin{definition}[Selection Rule]
A \emph{selection rule} $\mathcal{S}: \mathcal{E} \to \{1, \ldots, n\}$
maps the current input context $e$---including prompt structure,
register, distributional features, and semantic similarity to prior
acquisition contexts---to an index $k$, activating frame $F_k$ as
the operative frame for generation.
\end{definition}

\begin{definition}[Selection Stability Capacity (SSC)]
Let a \emph{task specification} $T$ determine a target frame $F^*$
as the frame whose output function best approximates the correct
response function for $T$ (independently of whether the model
produces hallucinated output). A model exhibits
\emph{Selection Stability Capacity} with respect to $F^*$ and $T$
if the selection rule $\mathcal{S}$ assigns $F^*$ across the full
distribution of inputs consistent with $T$. SSC failure occurs when
$\mathcal{S}$ activates a frame $F_k \neq F^*$ for some input $e$
in that distribution, producing output coherent within $F_k$ but
incorrect with respect to $T$.

\emph{Note on non-circularity.} The target frame $F^*$ is fixed by
the task specification $T$, not by whether the output is hallucinated.
Hallucination is then defined as SSC failure---output generated under
$F_k \neq F^*$---rather than the reverse. This ordering breaks the
apparent circularity between ``wrong frame'' and ``hallucinated output'':
the frame is wrong because it mismatches $T$, not because it produces
hallucination.
\end{definition}

\begin{remark}[Relation to prior work and Transformer evidence]
The interpretive stack model is developed formally in
\cite{minamoto2026accumulation}. The present paper
applies this model to hallucination without assuming familiarity with
the full formal apparatus.

The claim that Transformers implement something like frame selection
is supported by existing mechanistic interpretability findings.
Induction heads \cite{olsson2022induction} implement a form of
context-sensitive pattern completion that behaves differently
depending on the distributional context of the input---a primitive
form of frame selection. Attention head specialization
\cite{elhage2021mathematical} shows that different heads implement
qualitatively distinct computational roles, and that these roles
can be causally isolated via activation patching. Sparse
autoencoder decompositions \cite{bricken2023monosemanticity} have
identified monosemantic features corresponding to specific concepts,
domains, and registers within a single model's activation space.

The selection rule $\mathcal{S}$ is therefore not a posited
abstraction but a functional description of the mechanism by which
context-sensitive attention patterns route computation through
different circuit subgraphs. Its formal characterization in terms
of specific architectural components is the subject of OP2.
\end{remark}

\subsection{The Selection Rule}
\label{sec:selection-rule}

The frame that is active at any moment is determined by a
\emph{selection rule} $\mathcal{S}$.

\begin{definition}[Selection Rule]
\[
\mathcal{S}(e)
  = \operatorname{arg\,max}_k\;
    \mathrm{softmax}_k\!\left(\frac{s(e,\, C_k)}{\tau}\right)
\]
where $e \in \mathcal{E}$ is the current input context,
$C_k$ is the acquisition context of frame $F_k$,
$s(\cdot,\cdot): \mathcal{E} \times \mathcal{E} \to [0,1]$ is a context-similarity function,
and $\tau > 0$ is a temperature parameter.
\end{definition}

\begin{remark}[Cognitive interpretation of $\tau$]
The temperature $\tau$ operationalizes context-discrimination capacity.
Low $\tau$ produces sharp selection (dominant frame wins reliably);
high $\tau$ flattens the distribution, increasing misactivation risk.
In LLMs, $\tau$ effectively rises when input register diverges from
the fine-tuning distribution, when context is long, or when
chain-of-thought scaffolding is absent.
Each of these maps directly to a testable hallucination trigger
in Predictions~1--3: register shift (Type II), context length
(Type III), and prompt formulation (Type V).
\end{remark}

\noindent
Under this rule, the hallucination mechanism is precise:
when $e$ resembles the acquisition context $C_k$ of an archived
wrong-frame $F_k$ more closely than it resembles $C_1$ (the target
frame), $\mathcal{S}$ selects $F_k$, producing hallucination.
The temperature $\tau$ models context-discrimination capacity:
high $\tau$ (low discrimination) increases reactivation risk;
low $\tau$ (sharp discrimination) stabilizes the target frame.
This connects directly to Prediction~2 (Context-Structured
Hallucination Rate): hallucination rate should increase
monotonically with $s(e, C_k)$ for an archived wrong frame.

\subsection{Constraints on the Frame Concept}
\label{sec:frame-constraints}

A potential objection is that the frame concept is too expressive:
if any input--output pattern can be attributed to ``some frame,''
the account explains everything and therefore predicts nothing.
We address this by specifying three constraints that make the frame
concept falsifiable.

\textbf{Constraint 1 (Coherence).}
A frame $F$ must generate internally coherent outputs across its
activation domain: if inputs $e_1$ and $e_2$ activate the same
frame, their outputs should be stylistically, registrally, and
semantically consistent with each other and with an identifiable
training context.
This rules out treating every individual output as a distinct frame.

\textbf{Constraint 2 (Activation patching discriminability).}
Two proposed frames $F_1$ and $F_2$ are empirically distinct only
if there exists a pair of inputs $(e_1, e_2)$ such that activation
patching from $e_1$ to $e_2$ produces a \emph{qualitative} shift
in output type---not merely a quantitative change in token
probabilities.
This gives a necessary condition for frame distinctness that is
testable with existing mechanistic interpretability tools
\cite{conmy2023towards}.

\textbf{Constraint 3 (Acquisition context specificity).}
Each frame must be associated with an acquisition context $C$---a
characterizable region of the training distribution---such that
inputs with high distributional similarity to $C$ preferentially
activate the frame.
Distributional similarity can be operationalized via any of the
following: embedding proximity in the model's residual stream
at an early layer, prompt-template surface overlap (vocabulary,
register, formatting), or estimated fine-tuning corpus overlap
(when fine-tuning data is available).

The prediction holds under any proxy satisfying the following
\emph{pairwise ordering condition}: for any two inputs $e_1, e_2$
and target frame $F_k$ with acquisition context $C_k$,
\[
\delta(e_1, C_k) > \delta(e_2, C_k)
\;\Rightarrow\;
P(\mathcal{S}(e_1) = k) \geq P(\mathcal{S}(e_2) = k).
\]
This is a rank-preservation requirement, not a functional form
requirement: the prediction is that the Spearman rank correlation
between $\delta(\cdot, C_k)$ and the reactivation rate
$P(\mathcal{S}(\cdot) = k)$, measured over a test set of inputs,
is strictly positive.
This formulation is proxy-agnostic and pre-registerable.
This constraint generates the falsifiable prediction of
Prediction~2 (Section~\ref{sec:predictions}): reactivation rate
should be monotonically related to input similarity to $C$,
not uniformly distributed.

Together, these constraints rule out the degenerate case where
``frame'' is just a label attached post-hoc to any observed output
pattern.
The taxonomy in Section~3 instantiates each type with a specific
acquisition context, making the proposed frames falsifiable through
Constraints~1--3.

%-------------------------------------------------------------------
\section{A Frame-Selection Taxonomy of Hallucination}
%-------------------------------------------------------------------

We propose that hallucination types can be classified by the
mechanism of frame-selection failure. We identify five classes.

\begin{remark}[Relation to existing hallucination taxonomies]
Existing surveys \cite{huang2023survey, zhang2023siren} classify
hallucinations primarily by their output manifestation (factual,
faithfulness, contextual, etc.) and by attributed cause (knowledge
gap, attention failure, decoding artifact). The present taxonomy
organizes hallucinations by their \emph{frame-selection failure
mechanism}, which cuts across manifestation-based categories. A
factual hallucination may be a Type~I failure (distributional frame
bleed) or a Type~III failure (context boundary failure); a contextual
hallucination may be a Type~II failure (register-induced frame switch)
or a Type~IV failure (narrative frame overflow). The present taxonomy
is therefore not a replacement for manifestation-based taxonomies
but a proposed mechanistic layer beneath them, with distinct
implications for debiasing targets.
\end{remark}

\subsection{Type I: Distributional Frame Bleed}

The model is presented with an input that shares surface features
with a high-frequency training context, and the selection rule
activates the frame associated with that training context rather than
the intended frame.

\textbf{Example.} A question about a rare historical figure activates
the ``famous biography'' frame, generating plausible-sounding but
false biographical details that fit the statistical pattern of
well-documented figures.

\textbf{Mechanism.} The selection rule $\mathcal{S}$ uses surface
similarity to acquisition contexts as its primary signal. When the
input resembles a high-frequency training context more than the
low-frequency truth, the wrong frame wins.

\textbf{Calibration account.} The calibration account attributes
this to overconfidence in a low-evidence region. The frame-selection
account makes a stronger structural claim: even if the model were
perfectly calibrated, it would generate the wrong content as long
as the wrong frame is active, because the frame determines what
counts as coherent output.

\subsection{Type II: Register-Induced Frame Switch}

A shift in input register (formality, domain vocabulary, persona
cue) triggers a frame transition that is inappropriate for the
task.

\textbf{Example.} A model instructed to act as a strict factual
assistant generates speculative content when the user shifts to
a conversational register, because the conversational register
activates a ``chatbot'' frame with different generation norms.

\textbf{Mechanism.} The selection rule is sensitive to stylistic
features of the input. A register shift that resembles the
acquisition context of a different frame can override the intended
frame even when the task instruction remains explicit.

\subsection{Type III: Context Boundary Failure}

The model generates content from a document or context not present
in the current input window, treating the current context as
continuous with a related prior context.

\textbf{Example.} A model asked to summarize document A introduces
claims from document B, which is distributionally adjacent to
document A in the model's training data and shares surface features
with the current input, causing the selection rule to activate a
frame whose acquisition context encompasses document B's domain.

\textbf{Mechanism.} The selection rule activates a frame whose
acquisition context includes the adjacent document, effectively
extending the context boundary beyond the current input.

\subsection{Type IV: Narrative Overcompletion}

The model generates output that continues a narrative trajectory
consistent with its training distribution, overriding explicit
constraints in the prompt.

\textbf{Example.} A model asked to write a story with an unusual
ending instead generates a conventional ending that fits the
genre template, because the narrative completion frame is more
strongly activated than the ``constraint following'' frame.

\textbf{Mechanism.} Inputs with narratively structured context
produce selection outcomes in which the narrative-completion frame
is assigned by $\mathcal{S}$ with higher probability than the
constraint-following frame, because the narrative structure more
closely resembles the narrative frame's acquisition context.
The prompt constraint is processed within the already-active
narrative frame rather than overriding the selection outcome.

\subsection{Type V: Epistemic Frame Displacement}

The model responds as though it has access to information it does
not have, or as though it lacks information it does have, because
the active frame determines what counts as ``known.''

\textbf{Example.} A model confidently asserts outdated information
because the frame active for that domain was acquired during a
period when that information was correct, and the selection rule
does not update to a more recent frame.

\textbf{Mechanism.} The selection rule activates a temporally
anchored frame. Within that frame, the outdated information is
correct and confident assertion is appropriate. The error is not
overconfidence per se but temporal frame misalignment.

\textbf{Distinction from knowledge cutoff failure.}
Type~V is structurally distinct from ordinary knowledge cutoff
failure, in which the model lacks the relevant information
entirely.
The distinguishing test is frame-signal injection: present the
model with an explicit temporal marker (e.g., ``as of 2026'')
alongside the query.
Under a knowledge cutoff failure, the temporal marker cannot
override the error because the updated information is absent from
the weight space.
Under Type~V, the temporal marker functions as a frame-selection
signal, biasing the selection rule toward a more recent acquisition
context and suppressing the outdated frame---so the error
\emph{does} resolve when the frame signal is provided and
\emph{re-emerges} when it is removed.
This reactivation pattern under signal removal is the signature
prediction of the frame-selection account for Type~V, and is
not predicted by the knowledge cutoff account.

\emph{Concrete prompt pair.}
Consider a domain where the model's training data spans
two periods with contradictory factual content (e.g., a
leadership position that changed mid-training).
\begin{itemize}
\item \textbf{Baseline prompt:} ``Who is the CEO of [company]?''
  The model asserts the pre-change name confidently.
\item \textbf{Frame-signal prompt:} ``As of early 2026,
  who is the CEO of [company]?''
  Under Type~V, the temporal marker shifts the selection rule
  toward a later acquisition context; the model now returns
  the post-change name.
\item \textbf{Removal probe:} Remove the temporal marker and
  requery in the same conversation turn.
  Type~V predicts reversion to the earlier name; knowledge
  cutoff failure predicts no change (the model had accessed
  the correct information and should retain it in context).
\end{itemize}
The reversion in the removal probe is the critical observation:
it indicates that the correct information was available but
suppressed by frame reactivation, not that it was absent.

\subsection{Summary Table}

\begin{center}
\begin{tabular}{llll}
\toprule
Type & Name & Trigger & Mechanism \\
\midrule
I & Distributional Bleed & Surface similarity & High-frequency frame dominance \\
II & Register Switch & Stylistic shift & Register-frame coupling \\
III & Context Boundary & Adjacent context & Context window ambiguity \\
IV & Narrative Overcompletion & Narrative structure & Completion prior dominance \\
V & Epistemic Displacement & Temporal anchoring & Temporally fixed frame \\
\bottomrule
\end{tabular}
\end{center}

%-------------------------------------------------------------------
\section{Reinterpreting RAG, CoT, and System Prompts}
%-------------------------------------------------------------------

Under the frame-selection account, the major practical techniques
for reducing hallucination are not primarily knowledge augmentation
or reasoning support but \emph{interventions on the selection rule}.

\subsection{Retrieval-Augmented Generation (RAG)}

The standard interpretation of RAG is that it provides the model
with external knowledge that it lacks internally. The frame-selection
account offers a complementary interpretation: RAG modifies the
input context $e$ in ways that bias the selection rule toward
frames consistent with the retrieved content.

Specifically, retrieved documents carry distributional signals
(vocabulary, register, domain markers) that shift the similarity
between $e$ and frame acquisition contexts. By injecting content
that closely resembles the acquisition context of the target frame,
RAG increases the probability that $\mathcal{S}$ selects the
appropriate frame.

This interpretation makes a prediction that the knowledge
augmentation account does not: RAG should reduce hallucination
even when the retrieved content is semantically relevant but
factually redundant---because the frame-selection effect operates
through distributional similarity, not through logical entailment.

\subsection{Chain-of-Thought Prompting}

CoT prompting is standardly interpreted as enabling the model to
perform multi-step reasoning by externalizing intermediate steps.
The frame-selection account offers a structural reinterpretation:
CoT \emph{sequentially fixes} the interpretive frame at each
reasoning step, preventing frame drift between steps.

By making the reasoning process explicit in the context, CoT
continuously updates the input $e$ with the current reasoning
state, biasing the selection rule toward frames consistent with
the established reasoning trajectory. This reduces the probability
of a frame switch mid-reasoning that would generate output
inconsistent with the established conclusion.

The frame-selection account predicts that CoT should be especially
effective for Type IV (narrative overcompletion) and Type II
(register-induced frame switch) hallucinations, where the primary
risk is frame drift during an extended generation sequence.

\subsection{System Prompts and Persona Instructions}

System prompts function as \emph{prior frame injections}: by
establishing a role, persona, or task frame before any user input,
they bias the selection rule toward the specified frame for the
duration of the interaction.

However, the frame-selection account also explains the well-known
fragility of persona instructions under adversarial prompting
\cite{perez2022ignore, wei2023jailbroken}: a sufficiently strong
input signal---one that strongly resembles the acquisition context
of a different frame---can override the system prompt's prior
injection, activating the alternative frame and producing output
inconsistent with the intended persona.

This is not a failure of instruction following (the model
``understands'' the system prompt) but a failure of SSC: the
selection rule is not stable under the adversarial input.

%-------------------------------------------------------------------
\section{Distinguishing Predictions}
\label{sec:predictions}
%-------------------------------------------------------------------

The frame-selection account and the calibration account make
distinct empirical predictions. We identify four predictions
that are tractable with current mechanistic interpretability methods.

The following table maps the five hallucination types (Section~3)
to the predictions that most directly test each type's proposed
mechanism.

\begin{center}
\small
\begin{tabular}{lll}
\toprule
Hallucination Type & Primary Prediction & Secondary Prediction \\
\midrule
I: Distributional Bleed & P1 (Frame-coherent errors) & P2 (Context-structured rate) \\
II: Register Switch & P2 (Context-structured rate) & P4 (Circuit identity) \\
III: Context Boundary & P3 (RAG frame signal) & P4 (Circuit identity) \\
IV: Narrative Overcompletion & P1 (Frame-coherent errors) & P4 (Circuit identity) \\
V: Epistemic Displacement & P2 (Context-structured rate) & P4 (Circuit identity) \\
\bottomrule
\end{tabular}
\end{center}

\medskip

\noindent\textbf{Additional predictions for Types I--IV.}
Beyond the four primary predictions below, we note one
confirmable prediction per hallucination type that is tractable
with existing datasets:
\textit{Type I} (Distributional Bleed) predicts that clustered
hallucinated outputs will show lower perplexity \emph{within}
the misactivated frame than the base model's perplexity on the same outputs;
\textit{Type II} (Register Switch) predicts that hallucination rate
increases when system-prompt register mismatches the query register
(e.g., formal system prompt + casual query);
\textit{Type III} (Context Boundary) predicts that RAG retrieval
of frame-congruent but factually redundant content reduces hallucination
more than factually novel but frame-incongruent content;
\textit{Type IV} (Narrative Overcompletion) predicts that hallucinated
continuations show higher cosine similarity to the narrative frame's
training distribution than to the factual frame's distribution.

\subsection{Prediction 1: Frame-Coherent Error Structure}

\textbf{Calibration account (strongest version).} The most
sophisticated calibration accounts \cite{kadavath2022, xiong2024}
do not claim that hallucinated content is uniformly distributed.
They allow that hallucinated content is biased toward high-frequency
patterns in the training distribution (since overconfident regions
are precisely the high-frequency regions). On the strongest version
of the calibration account, hallucinated content should correlate
with training frequency and should be detectable by entropy-based
uncertainty measures.

\textbf{Frame-selection account.} Hallucinated content should
be coherent within the misactivated frame: factually wrong
but stylistically, registrally, and semantically appropriate to
the wrong frame. Critically, hallucinated content should be
\emph{low-entropy within the wrong frame}---it should appear
confident and internally consistent rather than uncertain---
because the model is generating correctly within an incorrect frame,
not generating incorrectly within a correct frame. Errors should
cluster by frame type rather than by training frequency.

\textbf{Test.} Cluster hallucinated outputs by semantic and
stylistic similarity. The frame-selection account predicts
cluster structure that corresponds to identifiable training
frames (biography frame, narrative frame, domain expert frame,
etc.).

\emph{Measurement order.}
To avoid circularity (Section~\ref{sec:frame-constraints},
Constraint~2), the test proceeds in two stages.
First, candidate frames are identified using the taxonomy of
Section~3 as a prior: each hallucination type specifies a
candidate acquisition context from which a frame hypothesis
can be derived (e.g., ``this Type~I error should cluster with
the high-frequency biography training context'').
Second, activation patching is used to confirm or disconfirm
the proposed frame assignment: if patching from a confirmed
biography-frame input to the hallucinated output leaves output
style qualitatively unchanged, the frame assignment is supported.
This two-stage procedure grounds the clustering test in
Constraint~2 without presupposing the result.

\subsection{Prediction 2: Context-Structured Hallucination Rate}

\textbf{Calibration account (strongest version).} Sophisticated
calibration accounts \cite{lin2022, xiong2024} predict that
hallucination rate should vary primarily with the model's
epistemic uncertainty about the topic, as estimated by consistency
across multiple samples or by probing-based uncertainty measures.
On this view, varying input register or vocabulary while holding
topic constant should not substantially affect hallucination rate
if topic-level uncertainty is controlled---the model's epistemic
state about a topic should be relatively invariant to how the
question is phrased.

\textbf{Frame-selection account.} Hallucination rate should vary
with the \emph{distributional distance} between the input context
and the target frame's acquisition context. Inputs that
distributionally resemble a high-frequency alternative frame
should hallucinate more than inputs that do not, independent of
topic coverage. This prediction is \emph{orthogonal} to
uncertainty: a topic the model ``knows well'' (low entropy) can
still hallucinate if register shift activates the wrong frame.

\textbf{Test.} Systematically vary input register and vocabulary
while holding topic constant. The frame-selection account predicts
that inputs resembling alternative frame acquisition contexts
will hallucinate more, even for topics with high coverage and
low model uncertainty.

\subsection{Prediction 3: RAG Effectiveness via Frame Signal}

\textbf{Knowledge augmentation account.} RAG should be effective
primarily when retrieved content provides factual information
the model lacks.

\textbf{Frame-selection account.} RAG should be effective
even when retrieved content is factually redundant, if it
provides strong frame-selection signals (vocabulary, register,
domain markers). Conversely, RAG should fail when retrieved
content provides correct facts but weak frame signals.

\textbf{Test.} Compare RAG conditions in which retrieved content
is (a) factually novel and frame-congruent, (b) factually
redundant and frame-congruent, (c) factually novel but
frame-incongruent. The frame-selection account predicts that
condition (b) reduces hallucination more than condition (c).

\subsection{Prediction 4: Circuit Persistence of Misactivated Frames}

\textbf{Calibration account.} Hallucination-prone regions of
the input space should correspond to regions of high model
uncertainty, identifiable by entropy in output distributions.

\textbf{Frame-selection account.} Hallucination-prone inputs
should activate identifiable circuits associated with the
misactivated frame, detectable via activation patching or
causal tracing \cite{elhage2021mathematical, conmy2023towards}.
The hallucinating model should show high activation of
wrong-frame circuits and low activation of correct-frame
circuits.

\textbf{Test.} Apply mechanistic interpretability methods to
compare circuit activation patterns on hallucinating vs.\
non-hallucinating inputs. The frame-selection account predicts
that the difference is in circuit identity (which frame is
active) rather than in output entropy alone.

\subsection{Negative Prediction: Falsification Condition}

To constitute a scientific theory rather than a post-hoc
classification scheme, the frame-selection account must
specify conditions under which it is falsified.
We state one critical negative prediction:

\textbf{Negative Prediction (Type~V).}
If the Type~V (epistemic frame displacement) account is correct,
then presenting an explicit temporal marker (``as of [date]'')
alongside a query where the model asserts outdated information
should \emph{reduce} the assertion rate, and \emph{removing}
the temporal marker from context should restore the original
assertion rate.
\emph{If neither the reduction nor the restoration occurs},
the frame-selection account of Type~V is falsified for that
instance, and the failure is better attributed to knowledge
cutoff (absence of information) rather than frame misactivation
(information present but suppressed).

The critical test is the restoration step: knowledge cutoff
failure predicts no change upon marker removal (the information
is simply absent from weight space); frame misactivation
predicts restoration (the suppressed frame reactivates when
the frame signal is removed).

\subsection{Competing Hypotheses and Their Predictions}

Two alternative explanations merit explicit comparison:

\textbf{Recency bias hypothesis.}
Outputs reflect recent context tokens more than semantically
appropriate earlier tokens, due to attention decay over distance.
On this view, hallucination follows from recency weighting,
not frame selection.
\emph{Distinguishing prediction:} recency bias predicts that
hallucination rate should correlate with the distance of
relevant context tokens from the generation point, independent
of distributional similarity to alternative acquisition contexts.
The frame-selection account predicts the opposite: hallucination
should correlate with distributional similarity to alternative
frames, even when task-relevant context is at the same position.

\textbf{Context overwrite hypothesis.}
Instructions or retrieved content directly overwrite the model's
prior tendency, and failures arise when overwriting is incomplete.
On this view, stronger instructions should produce monotonically
lower hallucination rates.
\emph{Distinguishing prediction:} context overwrite predicts
that longer, more emphatic system prompts should consistently
reduce hallucination.
The frame-selection account predicts that prompt emphasis is
insufficient when the input context strongly matches an
alternative frame's acquisition context---in those cases,
frame-signal injection (distributional match) should dominate
instruction strength.
The removal probe test (Prediction~4 above) directly
distinguishes the two: context overwrite predicts that
removal of a frame signal should not restore prior behavior
(since the instruction remains), while frame misactivation
predicts restoration.

\subsection{Type V Quantification}

\begin{table}[ht]
\centering
\small
\begin{tabular}{lcccc}
\hline
Condition & $N$ prompts & Assertion rate & vs. baseline & Predicted by \\
\hline
Baseline ($e_0$) & 100 & $r_0$ & --- & All accounts \\
Frame signal ($e_C$) & 100 & $r_C > r_0$ & $+\delta_1$ & Frame-sel., Bayesian \\
Signal removed ($e_{\text{rm}}$) & 100 & $r_{\text{rm}} \approx r_0$ & $-\delta_2$ & Frame-sel. only \\
\hline
\end{tabular}
\caption{Type~V (epistemic frame displacement) removal probe protocol.
$r_0$: baseline assertion rate of outdated information;
$r_C$: assertion rate under temporal-context frame signal;
$r_{\text{rm}}$: assertion rate after signal removal.
The frame-selection account predicts $r_C - r_{\text{rm}} > \delta_2$
(restoration upon removal); the knowledge cutoff account predicts
$r_C - r_{\text{rm}} \approx 0$.
Pre-registration requires specifying $\delta_1, \delta_2 > 0$
as multiples of within-condition standard deviation before data
collection. $N = 100$ per condition yields $>$90\% power for
effect sizes $\geq 0.3$ SD at $\alpha = 0.05$.}
\label{tab:typev}
\end{table}

%-------------------------------------------------------------------
\section{Implications for Debiasing}
%-------------------------------------------------------------------

\subsection{Why Confidence Calibration Is Insufficient}

If hallucination is primarily a frame-selection failure, then
improving confidence calibration addresses the symptom but not
the cause. A perfectly calibrated model that activates the wrong
frame will generate output that is confidently and coherently
wrong within that frame. The calibration intervention would
make the model \emph{less certain} about its wrong output,
but it would not prevent the wrong frame from being activated.

The appropriate target for debiasing is SSC: the model's ability
to maintain the correct frame across the full distribution of
intended inputs, including inputs that distributionally resemble
alternative frames.

\subsection{SSC-Targeted Interventions}

SSC can be targeted by interventions that modify the selection rule
rather than the content of individual frames. We identify three
candidate intervention types.

\textbf{Acquisition context poisoning.} During fine-tuning for a
target frame $F^*$, deliberately present training examples in
contexts that distributionally resemble alternative frame
acquisition contexts. This trains the selection rule to select
$F^*$ even in the presence of competing frame signals.

\textbf{Contrastive frame training.} Train the model to
explicitly distinguish frame-appropriate from frame-inappropriate
outputs for a given input, rather than only training on
frame-appropriate outputs. This may strengthen the selection
rule's sensitivity to frame mismatch.

\textbf{Context anchoring.} Modify inference-time prompting to
continuously re-anchor the frame signal, reducing the probability
of frame drift over long generation sequences. This is a
formalization of techniques already used in practice (e.g.,
repeating the system prompt at regular intervals).

\subsection{Implications for RAG and CoT Design}

Under the frame-selection account, RAG and CoT systems should
be designed not only to provide relevant content or reasoning
scaffolds but to maximize frame-selection signal:

\begin{enumerate}[label=(\roman*)]
\item Retrieved documents should be filtered not only for
semantic relevance but for frame congruence---documents that
strongly resemble the target frame's acquisition context.
\item CoT prompts should be designed to explicitly re-establish
the frame at each reasoning step, not only to provide
intermediate conclusions.
\item System prompts should be tested for SSC robustness:
how strongly must an adversarial input deviate from the system
prompt's frame signal before the model switches frames?
\end{enumerate}

%-------------------------------------------------------------------
\section{Relation to Prior Work}
%-------------------------------------------------------------------

The calibration account of hallucination has been developed
extensively \cite{kadavath2022, lin2022, xiong2024}. The
frame-selection account is complementary rather than competing:
we argue that calibration failure and frame misactivation are
distinct failure modes that coexist in practice, and that
current research over-indexes on calibration at the expense of
frame-selection analysis.

The concept of ``competing experts'' or ``mixture of behaviors''
in neural networks is not new \cite{jacobs1991, jordan1994},
and recent mechanistic interpretability research has identified
competing circuits in LLMs \cite{conmy2023towards, bricken2023monosemanticity}. Our contribution
is to organize these observations within a formal framework
(the interpretive stack) that connects them to hallucination
taxonomy and generates concrete testable predictions.

The gate function in mixture-of-experts (MoE) architectures
\cite{jacobs1991, jordan1994} implements a selection rule in
mathematically explicit form: a learned function assigns each input
to one or more expert modules. Frame misactivation corresponds,
under this analogy, to gating error. The interpretive stack model
differs from MoE in three respects. First, in standard MoE, experts
are trained jointly with the gate and are well-separated by design;
LLM frames are acquired sequentially across pretraining and fine-tuning
phases and may overlap substantially in input distribution. Second, MoE
gates are trained with explicit gating objectives; the LLM selection
rule $\mathcal{S}$ is emergent and not directly optimized.
Third, MoE models exhibit stable gating within their training
distribution by construction; LLM SSC failures occur specifically under
distributional shift relative to each frame's acquisition context.
The frame misactivation account therefore generalizes MoE gating
failure analysis to the sequential, multi-phase training setting
of large language models.

The taxonomy of hallucination types we propose differs from
existing taxonomies \cite{huang2023survey, zhang2023siren} in
its organizing principle: rather than classifying by output
type (factual vs.\ faithfulness hallucination, etc.), we
classify by the frame-selection mechanism responsible for the
error. This distinction matters because the two taxonomies
imply different debiasing targets.

The relation between fine-tuning and frame persistence is
developed in \cite{minamoto2026accumulation}, where fine-tuning is
shown to implement frame accumulation rather than overwriting,
with prior frames remaining selectable under appropriate input
conditions. The present paper focuses on inference-time frame
selection rather than training-time frame accumulation.

A complementary perspective comes from thinking mid-training
\cite{lanchantin2026midtraining}, which introduces an intermediate
SFT+RL training phase between pretraining and post-training. The
finding that this improves post-training performance by 3.2$\times$
is consistent with the frame-selection account: models that develop
structured frame-transition patterns during mid-training exhibit
higher Selection Stability Capacity during post-training, and are
therefore less prone to the frame-misactivation failures analyzed
here. The training pipeline determines not only which frames are
acquired but how robustly the selection rule can discriminate among
them under distributional shift.

%-------------------------------------------------------------------
\section{Pilot Empirical Evidence}
\label{sec:pilot}
%-------------------------------------------------------------------

The predictions of Section~\ref{sec:predictions} require
mechanistic interpretability experiments at a scale beyond
the scope of this theoretical paper.
We report a pilot study using representational similarity
analysis (RSA) that provides initial indirect support for
the frame-selection account's central structural claim:
that training establishes a stable relational geometry that
is sensitive to distributional input context.

\paragraph{Setup.}
We applied RSA to three training checkpoints of Pythia-70m
\cite{biderman2023}: step~0 (random initialization),
step~1000 (early training), and step~143000 (fully trained
reference).
Pairwise dissimilarity matrices were computed over $n = 29$
probe concepts at hidden layers $\ell \in \{2, 4, 6\}$,
using three prompt templates averaged to reduce surface-form
bias.
Structural similarity was measured as the second-order RSA
correlation $\rho_\mathrm{cog}(A, B; \ell) =
\mathrm{Spearman}(\mathrm{upper}(D^{A,\ell}),
\mathrm{upper}(D^{B,\ell}))$.

\paragraph{Results.}
Two findings are relevant to the frame-selection account.

\emph{Finding 1: Training induces stable relational geometry.}
$\rho_\mathrm{cog}$ between early and fully trained checkpoints
increases monotonically across all layers
(L6: step-0 $\to$ step-1000: $0.017 \to 0.535$),
indicating that training progressively installs a stable
relational structure rather than reorganizing it randomly.
This is consistent with the frame-selection account's
premise that trained models maintain interpretive frames
as structured attractors in weight space.

\emph{Finding 2: Cross-architecture convergence.}
$\rho_\mathrm{cog}(\text{GPT-2 trained},
\text{Pythia-trained}) \approx 0.42$ versus
$\rho_\mathrm{cog}(\text{GPT-2 trained},
\text{Pythia-step-0}) \approx 0.10$ across all layers.
Two architecturally distinct models trained on different
corpora converge to similar relational geometry,
while random initialization shares no such structure.
This suggests that the relational geometry captured by RSA
is a property of training-induced frames rather than
architectural artifact---consistent with the frame-selection
account's prediction that hallucination patterns should
be frame-coherent (Prediction~1) and context-sensitive
(Prediction~2) rather than architecture-specific.

\paragraph{Interpretation and limitations.}
These results are consistent with but do not directly test
Predictions~1--4, which require activation-level evidence
of frame identity during hallucination.
Critically, the RSA method cannot distinguish whether
distributional shift in the input \emph{changes} which
frame is selected (the frame-selection account's prediction)
or merely adds noise to an otherwise fixed representation.
Direct tests via the Removal Probe protocol
(Section~\ref{sec:predictions}, Prediction~4;
Table~\ref{tab:typev}) remain as the primary empirical priority.
All code is available at
\url{https://github.com/minamominamoto/grounding-without-reality}.

%-------------------------------------------------------------------
\section{Open Problems}
%-------------------------------------------------------------------

\begin{enumerate}[label=\textbf{OP\arabic*.}, leftmargin=*, itemsep=4pt]

\item \textbf{Mechanistic identification of frames.}
    Develop a method for identifying the boundaries of
    interpretive frames in an LLM's weight space via sparse
    autoencoders \cite{bricken2023monosemanticity}, activation
    patching, or probing. This would enable empirical validation
    of the five hallucination types.

\item \textbf{Formal characterization of the selection rule.}
    Determine the computational form of $\mathcal{S}$ in
    transformer architectures. Which attention heads, residual
    stream directions, or MLP layers implement context-to-frame
    mapping?

\item \textbf{SSC measurement.}
    Develop an operational measure of SSC for LLMs: given a
    target frame and an adversarial input distribution, estimate
    the probability of frame misactivation as a function of
    input distributional distance from the target frame's
    acquisition context.

\item \textbf{RAG frame-signal hypothesis.}
    Test Prediction 3 empirically: does RAG reduce hallucination
    via frame-signal injection independent of factual content?
    A controlled experiment varying frame congruence while
    holding factual content constant would provide the critical
    evidence.

\item \textbf{CoT as sequential frame fixing.}
    Determine whether the benefit of CoT is primarily
    attributable to frame stabilization rather than (or in
    addition to) intermediate reasoning. An ablation study
    comparing structured CoT (which maintains frame signal)
    against unstructured scratchpad (which may not) would
    be informative.

\end{enumerate}

%-------------------------------------------------------------------
\section{Conclusion}
%-------------------------------------------------------------------

We have proposed a structural account of LLM hallucination as
frame misactivation: the selection rule governing which
interpretive frame is active in generation selects an
inappropriate frame for the current input context, producing
output that is fluent and coherent within the wrong frame.

This account does not replace the calibration account but
complements it: calibration failure and frame misactivation
are distinct failure modes that call for distinct interventions.
The appropriate intervention target for frame misactivation
is Selection Stability Capacity (SSC)---the model's ability
to maintain the correct frame across the intended input
distribution---rather than confidence calibration alone.

Under this account, RAG and CoT are reinterpreted as
frame-selection interventions, and their fragility under
adversarial inputs is explained as SSC failure rather than
instruction-following failure. The five hallucination types
we identify generate concrete, falsifiable predictions that
distinguish the frame-selection account from the calibration
account, providing an empirical program for validation via
mechanistic interpretability.

%-------------------------------------------------------------------
\section*{Acknowledgments}
%-------------------------------------------------------------------

No funding was received for this work.

%-------------------------------------------------------------------
\begin{thebibliography}{99}

\bibitem{minamoto2026accumulation}
Minamoto, M. (2026). Abductive accumulation: Why interpretive frames
persist, and how they drive hallucination in language models.
Preprint (not yet on arXiv). Available at:
\texttt{https://github.com/minamominamoto/topology-of-grounding}

\bibitem{kadavath2022}
Kadavath, S., et al. (2022). Language models (mostly) know what
they know. \textit{arXiv:2207.05221}.

\bibitem{lin2022}
Lin, S., Hilton, J., \& Evans, O. (2022). TruthfulQA: Measuring
how models mimic human falsehoods. \textit{Proceedings of ACL 2022},
3214--3252.

\bibitem{xiong2024}
Xiong, M., et al. (2024). Can LLMs express their uncertainty?
An empirical evaluation of confidence elicitation in LLMs.
\textit{Proceedings of ICLR 2024}.

\bibitem{perez2022ignore}
Perez, F., \& Ribeiro, I. (2022). Ignore previous prompt:
Attack techniques for language models. \textit{arXiv:2211.09527}.

\bibitem{wei2023jailbroken}
Wei, A., Haghtalab, N., \& Steinhardt, J. (2023). Jailbroken:
How does LLM safety training fail? \textit{NeurIPS 2023}.

\bibitem{elhage2021mathematical}
Elhage, N., et al. (2021). A mathematical framework for transformer
circuits. \textit{Transformer Circuits Thread}.
\texttt{https://transformer-circuits.pub/2021/framework/index.html}

\bibitem{olsson2022induction}
Olsson, C., et al. (2022). In-context learning and induction heads.
\textit{Transformer Circuits Thread}.
\texttt{https://transformer-circuits.pub/2022/in-context-learning-and-induction-heads/index.html}

\bibitem{conmy2023towards}
Conmy, A., Mavor-Parker, A., Lynch, A., Heimersheim, S.,
\& Garriga-Alonso, A. (2023). Towards automated circuit discovery
for mechanistic interpretability. \textit{NeurIPS 2023}.

\bibitem{bricken2023monosemanticity}
Bricken, T., et al. (2023). Towards monosemanticity: Decomposing
language models with dictionary learning.
\textit{Transformer Circuits Thread}.

\bibitem{lanchantin2026midtraining}
Lanchantin, J., Li, D., Nguyen, T., Sukhbaatar, S., Kulikov, I.,
Weston, J., \& Li, X. (2026). Thinking mid-training: Reinforcement
learning of interleaved reasoning. In Tan, E. X., et al.,
\textit{Self-Improving Pretraining}. \textit{arXiv:2601.21343}.

\bibitem{jacobs1991}
Jacobs, R. A., Jordan, M. I., Nowlan, S. J., \& Hinton, G. E.
(1991). Adaptive mixtures of local experts. \textit{Neural
Computation}, 3(1), 79--87.

\bibitem{jordan1994}
Jordan, M. I., \& Jacobs, R. A. (1994). Hierarchical mixtures
of experts and the EM algorithm. \textit{Neural Computation},
6(2), 181--214.

\bibitem{huang2023survey}
Huang, L., et al. (2023). A survey on hallucination in large
language models: Principles, taxonomy, challenges, and open
questions. \textit{arXiv:2311.05232}.

\bibitem{zhang2023siren}
Zhang, Y., et al. (2023). Siren's song in the AI ocean: A survey
on hallucination in large language models.
\textit{arXiv:2309.01219}.

\bibitem{biderman2023}
Biderman, S., et al.\ (2023).
Pythia: A suite for analyzing large language models across training and scaling.
\textit{Proceedings of ICML 2023}.

\end{thebibliography}

\end{document}
